\section{Project Strategy \& Scope}

\subsection{Mission Statement}
The primary objective of \textbf{Project Pienza} is to model and optimize the decision-making policy of an expert agent within the ride-hailing ecosystem. The environment is formally defined as a \textit{sequential, non-cooperative game} governed by intrinsic information asymmetry. 

Unlike traditional "top-down" studies relying on aggregated data, Pienza employs a \textbf{"Bottom-Up"} methodology grounded in the operational reality of Mexico City. Through the longitudinal quantification of a single Expert Agent ($N=1$), this research models the micro-economic physics of the Gig Economy.

\subsection{Scope \& Constraints}
The project explicitly rejects the premise of reverse-engineering the proprietary pricing algorithms of ride-hailing platforms. Modeling a dynamic, global algorithm using a local boutique dataset is statistically unfeasible. 

Instead, Pienza reorients the analytical lens—shifting away from the fare prediction models common in public repositories—to focus on the sole variable under absolute control: \textbf{The Agent's Decision}. By pivoting to Behavioral Economics and Decision Science, the project transforms a complex prediction task into a solvable policy optimization endeavor.

\subsection{Core Framework}
The project is architected across three strategic pillars that bridge the gap between tactical fieldwork and scalable autonomous intelligence:

\begin{enumerate} 
    \item \textbf{High-Fidelity Environment Digitization (The Foundation):} 
    To overcome the \textit{Survivorship Bias} inherent in ride-hailing data, a custom dual-engine ingestion pipeline was deployed to capture the \textbf{Total Offer Stream}. By utilizing OCR-based archival and geospatial triangulation, the system digitizes the "Negative Class" (rejected offers), providing the necessary ground truth to map the agent's exact decision boundary.
    
    \item \textbf{Hierarchical Policy Cloning (The Engine):} 
    The core analytical breakthrough involves the transition from monolithic modeling to a \textbf{Cognitive Cascade}. This hierarchical inference architecture utilizes XGBoost to replicate the agent's multi-tiered decision logic—separating deterministic operational "triage" (geospatial/economic filters) from nuanced opportunity-cost "gambles". 
    
\item \textbf{Generative Synthesis \& Topological Scaling (The Frontier):} 
    The final pillar transforms the single-agent study into a scale-invariant simulation. Using \textbf{Conditional GANs (cGAN)}, the project synthesizes a 1,000,000-row data manifold that preserves the non-linear physics of the marketplace as encountered during the observation window.\footnote{The synthetic marketplace functions as an infinite Keras generator of the specific market regime and "lived experience" captured in the dataset, rather than a generalizable, dynamic model of total city-wide liquidity.} This environment serves as the sandbox for \textbf{Network Graph Analysis}, where the city is modeled as a \textit{Tensor of Potentials} to orchestrate autonomous fleet routing and maximize steady-state value.
    
\end{enumerate}

\subsection{Initial Hypothesis: The Payout Spread}
Research began with a longitudinal field observation regarding fare elasticity. The hypothesis stated that the system applied an implicit efficiency penalty, observed as a discrepancy between the promised fare and realized earnings. A preliminary pilot study ($N \approx 150$, July-August 2025 cohort) revealed a \textbf{Payout Rate (Spread)} between $75\%$ and $85\%$ of the Base Fare. The distribution followed session-based moving averages, suggesting intra-block temporal stability driven by market conditions.

An initial model correlated this Spread with the \textbf{Time Delta} (the discrepancy between estimated and realized trip time). The goal was to validate whether time savings generated a fare penalty.

\textbf{Key Finding:} A benchmark comparison showed that Simple Linear Regression outperformed base Machine Learning models. No significant non-linear relationship was identified; the correlation remained weak. While basic statistics resolved the regression problem for the pilot sample, scaling this to a robust ML model would require thousands of completed trips, representing an infeasible operational time cost for marginal analytical return.

\subsection{Strategic Pivot to Classification}
Based on logistical constraints and the low predictive value of the initial regression approach, the research objective was redefined:
\begin{itemize}
    \item \textbf{From:} Price variance prediction (Regression), which required scarce completed trip data.
    \item \textbf{To:} Behavior modeling (Classification), which utilized the total offer stream.
\end{itemize}

This shift resolved the data scarcity issue by enabling the inclusion of the \textbf{Negative Class} (rejected offers). By modeling decision policy rather than price, the project applies high-dimensional Machine Learning (XGBoost) to capture the agent's non-linear decision boundary.

\subsection{Project Roadmap}
The project followed an iterative development lifecycle where feature engineering evolved in response to the analytical findings of each stage:

\begin{description}[style=nextline, leftmargin=0.5cm, font=\bfseries]
    \item[Phase 1: Acquisition \& Ground Truth (Aug 22 -- Oct 1, 2025)]
    High-fidelity manual capture of operational data via the bespoke \textit{GTS Webapp}. Establishment of the target variable and the primary contextual ride attributes.

    \item[Phase 2: Data Engineering \& Architecture (Oct 2 -- Nov 20, 2025)]
    Transition to a relational database (\texttt{pienza.db}) and Star Schema design. Implementation of automated OCR pipelines and normalization protocols. Creation of the statefull \texttt{engineered\_features} table.

    \item[Phase 3: Exploratory Analysis \& Causal Inference (Nov 21 -- Dec 3, 2025)]
    Diagnostic audit of marketplace physics and rational search time boundaries. Causal modeling of baseline heteroscedasticity and quantification of the platform's inelastic 'Integrity Buffer.'

    \item[Phase 4: Unsupervised Learning \& Geo-Remediation (Dec 4, 2025 -- Jan 2, 2026)]
    Application of HDBSCAN for zone discovery and surgical coordinate cleaning using hand-crafted heuristic polygons. Generation of the \textbf{Silver Palette} (geo-semantic attributes).

    \item[Phase 5: Supervised Imitation Learning (Jan 3 -- Jan 12, 2026)]
    Model evaluation tournament and implementation of the \textit{Cognitive Cascade} hierarchical architecture for the champion model.

    \item[Phase 6: Generative Moonshots (Jan 13, 2026 -- March 31, 2026)\footnote{The extended duration of this phase reflects a non-linear, iterative workflow; it was executed concurrently with the formal authoring of this manuscript, culminating in the simultaneous completion of Phase 6 and the first comprehensive draft of this paper.}] 
    $O(1)$ neural spatial inference and cGAN manifold synthesis. Formulation of the internal Mobility Tensor to establish the Markov Decision Process (MDP) baseline for autonomous routing.
\end{description}

\subsection{A Note on System Architecture \& Evolution}
The rigorous architecture of the data ecosystem was established as a foundational objective. The infrastructure evolved through distinct operational states to match the complexity of the inquiry:

\begin{itemize}
    \item \textbf{Phase 1 (Acquisition Staging):} Data ingestion from Engines 1 \& 2 routed into \textit{Google Sheets} for initial validation.
    \item \textbf{Phase 2 (The Engineering Core):} Transition to a relational architecture. Execution of the idempotent ETL pipeline to forge the definitive \textbf{SQLite} asset (\texttt{pienza.db}).
    \item \textbf{Phases 3--5 (The Analytical Loop):} A file-based cloud architecture where Python Notebooks interacted with \texttt{pienza.db} as the immutable Single Source of Truth.
    \item \textbf{Phase 6 (Cloud Scaling):} Migration to \textbf{Google BigQuery} to support Generative AI workflows.
    \begin{itemize}
        \item \textbf{\texttt{pienza\_mini}:} A high-fidelity cloud replica of the ground-truth database, preserving the identical relational logic.
        \item \textbf{\texttt{pienza\_big}:} A 1,000,000-row synthetic dataset generated via GANs, persisted as Parquet artifacts and queried through BigQuery's external table interface.
    \end{itemize}
    \item \textbf{Interactive Deployment (Current State):} The project concludes with the transition from notebook-based exploration to an \textbf{Interactive White Paper} hosted on Streamlit. Leveraging \textit{GitHub Codespaces} and an \textit{orphan branch architecture}, this layer decouples the core research codebase from the operational interface, allowing for real-time model interrogation in a production-ready environment.
\end{itemize}